\documentclass[12pt,compress]{beamer}
\usepackage{ifthen,hyperref}

\title{How Muon Reconstruction is Done}
\author{Jim Pivarski}
\institute{Texas A\&M University}
\date{6 September, 2006}

\setbeamertemplate{navigation symbols}{}
\setbeamertemplate{headline}{\includegraphics[height=1 cm]{cmslogo} \hspace{0.1 cm} \includegraphics[height=1 cm]{tamulogo} \hfill
\begin{minipage}{9 cm}
\vspace{-0.75 cm} \small
\begin{center}
\ifthenelse{\equal{\insertpagenumber}{0}}{}{\insertsection}
\end{center}
\end{minipage} \hfill
\begin{minipage}{1 cm}
\vspace{-0.75 cm} \small
\begin{center}
\ifthenelse{\equal{\insertpagenumber}{0}}{}{\insertpagenumber/\pageref{numpages}}
\end{center}
\end{minipage}}

\xdefinecolor{verylightgray}{rgb}{0.95,0.95,0.95}
\beamertemplateshadingbackground{verylightgray}{white}

\begin{document}
\addtocounter{page}{-1}
\frame{\titlepage}
\section*{Muon Reconstruction --- Jim Pivarski}

\begin{frame}
\frametitle{The Procedure}
\begin{enumerate}
\item Translate raw data into Digis, Digis into RecHits
\item Resolve left/right ambiguity
\item Construct segments
\item Fit to a ``stand-alone'' muon track
\item Do all silicon tracking
\item Match silicon-only tracks with muon-only tracks
\item Combine the fits, present reco::Muon
\end{enumerate}

\vfill
\textcolor{darkgray}{Stand-alone muons are probably {\it not} used to seed silicon track-finding.}

%% \vfill
%% \begin{minipage}{\linewidth}
%% \scriptsize \textcolor{darkgray}{(There's no good reason to do it: if you
%% reconstruct a 12-layer muon track, there's no reason to confirm it in
%% the silicon before deciding to keep the event.  The situation is
%% different for electrons.)}
%% \end{minipage}

\end{frame}

\begin{frame}
\begin{center}
\includegraphics[width=0.8\linewidth]{muonreco.png}
\end{center}
\end{frame}

\begin{frame}
\frametitle{Local Reconstruction}
Drift tubes (DT) and Cathode Strip Chambers (CSC) are reconstructed independently.

\begin{center}
\includegraphics[width=0.7\linewidth]{local_reco.png}
\end{center}
\end{frame}

\begin{frame}
\frametitle{Local Reconstruction}

All the hits and segments are TrackingRecHits, the same object used in
silicon tracking.

\begin{center}
\includegraphics[width=\linewidth]{family_tree.pdf}
\end{center}
\end{frame}

\begin{frame}
\frametitle{Local Reconstruction Status}
Completed before June Annual Review

\begin{center}
\includegraphics[width=0.5\linewidth]{recolocalmuon.png}
\end{center}

10 May Tutorial

{\tt \scriptsize \href{https://uimfon.cern.ch/twiki/bin/view/CMS/MuTutorialSegments}{https://uimfon.cern.ch/twiki/bin/view/CMS/MuTutorialSegments}}
\end{frame}

\begin{frame}
\frametitle{Stand-alone Muon Reconstruction}
\begin{enumerate}
\item CosmicMuonSeedGenerator constructs TrajectorySeeds out of
matched segments (TrajectorySeed is the same object used by silicon
tracking)
\item StandAloneMuonProducer performs the fit and creates reco::Tracks
(same object as silicon tracking)
\end{enumerate}

\vfill
SteppingHelixPropagatorESProducer is the simplest track/errors
propagator through material and non-uniform magnetic field.  Two
alternatives are in development.

\vfill
9 May talk on propagation by Nicola Amapane

{\tt \scriptsize \href{http://indico.cern.ch/conferenceDisplay.py?confId=1929}{http://indico.cern.ch/conferenceDisplay.py?confId=1929}}

\end{frame}

\begin{frame}
\frametitle{Stand-alone Reconstruction Status}

6 July Tutorial: ``not yet a reliable tool.  We are working on it!''

{\tt \scriptsize \href{https://twiki.cern.ch/twiki/bin/view/CMS/July06MuonReco}{https://twiki.cern.ch/twiki/bin/view/CMS/July06MuonReco}}

\vfill
\begin{columns}
\column{0.5\linewidth}

22 August PRS/mu:

\vspace{0.25 cm}
``seems to be basically working in all muon detectors!''
\begin{center}
and \mbox{\hspace{2 cm}}
\end{center}
``formally working''

\column{0.5\linewidth}
\includegraphics[width=\linewidth]{recomuon.png}
\end{columns}
\end{frame}

\begin{frame}
\frametitle{Combining Muon Tracks with Silicon Tracks}

What I can determine from the code:
\begin{itemize}
\item GlobalMuonProducer inputs stand-alone tracks and outputs reco::Muons.
\item reco::Muon contains references to stand-alone muon track, silicon-only track, and combined track.
\item The actual work is done by these classes:
\end{itemize}

\begin{center}
\includegraphics[width=0.9\linewidth]{globaltrackfinder.png}
\end{center}
\end{frame}

\begin{frame}
\frametitle{Global Muon Fit Status}
\begin{center}
Probably a work in progress.
\end{center}
\end{frame}

\begin{frame}
\frametitle{Actively-pursued Alternative}

UCSB (Claudio Campagnari, Dmitriy Kovalskyi, Slava Krutelyov, Jacob
Ribnik) wish to simplify the process, replacing

\begin{center}
stand-alone muon track $+$ silicon-only track $\longrightarrow$ global track
\end{center}

with

\begin{center}
silicon-only track $+$ muon segments $\longrightarrow$ global track
\end{center}

\vfill See Dmitriy Kovalskyi's talk at 22 August PRS/mu meeting
{\tt \scriptsize \href{http://indico.cern.ch/conferenceDisplay.py?confId=5292}{http://indico.cern.ch/conferenceDisplay.py?confId=5292}}

\end{frame}

\begin{frame}
\frametitle{Where to find things}

\begin{itemize}\setlength{\itemsep}{0.25 cm}
\item There are {\it no} CMS-NOTES (except 2006/010 for ORCA)

\item The PRS/mu group meets every other Tuesday (most recently 5 September) 10:30--12:30 our time

{\tt \scriptsize \href{http://indico.cern.ch/categoryDisplay.py?categId=25}{http://indico.cern.ch/categoryDisplay.py?categId=25}}

\item June 2006 Annual Review

{\tt \scriptsize \href{http://indico.cern.ch/conferenceDisplay.py?confId=3247}{http://indico.cern.ch/conferenceDisplay.py?confId=3247}}

\item Muon Software TWiki ($\rightarrow$ tutorials)

{\tt \scriptsize \href{https://uimon.cern.ch/twiki/bin/view/CMS/MuonSW}{https://uimon.cern.ch/twiki/bin/view/CMS/MuonSW}}

\item Doxygen Muon Overview (only lists packages)

{\tt \scriptsize \href{https://uimon.cern.ch/twiki/bin/view/CMS/MuonSW}{https://uimon.cern.ch/twiki/bin/view/CMS/MuonSW}}

\end{itemize}
\label{numpages}
\end{frame}

\end{document}
